% =====================================================================
%  MANUEL D'UTILISATION --- Alanya
%  Application mobile de communication (Flutter / Android)
%
%  Compilation recommandee :
%     pdflatex (x2)  |  latexmk -pdf  |  Overleaf (pdfLaTeX)
%
%  Captures d'ecran : deposer les PNG dans le sous-dossier captures/.
%  Tant qu'une capture est absente, un cadre "Capture a inserer" s'affiche
%  a sa place (le document compile quand meme).
%       captures/01-login.png, captures/02-signup.png, ...
%  Logos : images/logo-enspy.png , images/logo-Alanya.png
% =====================================================================
\documentclass[11pt,a4paper,oneside]{report}

% --------------------------------------------------------------------
%  Encodage, polices, langue
% --------------------------------------------------------------------
\usepackage[utf8]{inputenc}
\usepackage[T1]{fontenc}
\usepackage{textcomp}
\usepackage{lmodern}
\usepackage[french]{babel}
\usepackage{microtype}
\usepackage{amssymb}     % \checkmark

% --------------------------------------------------------------------
%  Mise en page
% --------------------------------------------------------------------
\usepackage[a4paper,margin=2.4cm,headheight=15pt]{geometry}
\usepackage{graphicx}
\usepackage{xcolor}
\usepackage{array}
\usepackage{booktabs}
\usepackage{tabularx}
\usepackage{enumitem}
\usepackage{caption}
\usepackage{float}
\usepackage{fancyhdr}
\usepackage{titlesec}
\usepackage[skins]{tcolorbox}

% --------------------------------------------------------------------
%  Schemas
% --------------------------------------------------------------------
\usepackage{tikz}
\usetikzlibrary{positioning,arrows.meta,shapes.geometric,fit,calc,backgrounds}

% --------------------------------------------------------------------
%  Liens (charge en dernier)
% --------------------------------------------------------------------
\usepackage[hidelinks]{hyperref}

% --------------------------------------------------------------------
%  Palette de marque Alanya
% --------------------------------------------------------------------
\definecolor{AlanyaIndigo}{HTML}{3F51B5}
\definecolor{AlanyaIndigoDark}{HTML}{1A237E}
\definecolor{AlanyaIndigoStrong}{HTML}{303F9F}
\definecolor{AlanyaContainer}{HTML}{E8EAF6}
\definecolor{AlanyaInk}{HTML}{1A1D23}
\definecolor{AlanyaGray}{HTML}{5B6273}
\definecolor{AlanyaLine}{HTML}{CBD0E0}
\definecolor{AlanyaGreen}{HTML}{1FA363}
\definecolor{AlanyaAmber}{HTML}{B7791F}
\definecolor{AlanyaBlue}{HTML}{2E90FA}
\definecolor{boxBg}{HTML}{F6F7FB}

\hypersetup{
  colorlinks=true,
  linkcolor=AlanyaIndigoStrong,
  urlcolor=AlanyaIndigoStrong,
  pdftitle={Manuel d'utilisation - Alanya},
  pdfauthor={Equipe Alanya},
  pdfsubject={Guide utilisateur de l'application mobile Alanya}
}

% --------------------------------------------------------------------
%  En-tetes / pieds de page
% --------------------------------------------------------------------
\pagestyle{fancy}
\fancyhf{}
\fancyhead[L]{\small\nouppercase{\leftmark}}
\fancyhead[R]{\small\textcolor{AlanyaIndigo}{\textbf{Alanya}} --- Manuel d'utilisation}
\fancyfoot[C]{\small\thepage}
\renewcommand{\headrulewidth}{0.4pt}
\renewcommand{\footrulewidth}{0.0pt}
\fancypagestyle{plain}{%
  \fancyhf{}\fancyfoot[C]{\small\thepage}%
  \renewcommand{\headrulewidth}{0pt}}

% --------------------------------------------------------------------
%  Titres de chapitres / sections
% --------------------------------------------------------------------
\titleformat{\chapter}[display]
  {\normalfont\sffamily\bfseries\color{AlanyaIndigo}}
  {\filright\large CHAPITRE\ \Huge\thechapter}
  {0.6ex}{\Huge}
\titlespacing*{\chapter}{0pt}{10pt}{30pt}
\titleformat{\section}
  {\normalfont\sffamily\Large\bfseries\color{AlanyaIndigoDark}}{\thesection}{1em}{}
\titleformat{\subsection}
  {\normalfont\sffamily\large\bfseries\color{AlanyaIndigoStrong}}{\thesubsection}{1em}{}

% --------------------------------------------------------------------
%  Macros utilitaires
% --------------------------------------------------------------------
% Libelle de bouton / commande de l'interface
\newcommand{\bouton}[1]{\textbf{\textsf{#1}}}
% Libelle de champ de saisie / onglet (entre guillemets)
\newcommand{\champ}[1]{\og{}#1\fg{}}
% Titre de colonne de tableau
\newcommand{\thd}[1]{\textbf{\textcolor{AlanyaIndigoDark}{#1}}}

% Cadre de remplacement generique si image absente
\newcommand{\imageOuCadre}[3]{%
  \IfFileExists{#1}{\includegraphics[width=#2]{#1}}{%
    \fbox{\begin{minipage}[c][4.2cm][c]{#2}%
      \centering\sffamily\color{AlanyaGray}%
      \textbf{[ Image manquante ]}\\[6pt]\small #3%
    \end{minipage}}}%
}

% Capture d'ecran (format portrait telephone) numerotee automatiquement
\newcommand{\capture}[2]{%
  \begin{figure}[H]\centering
    \IfFileExists{captures/#1}{\includegraphics[width=6cm]{captures/#1}}{%
      \fbox{\begin{minipage}[c][8.5cm][c]{5.6cm}%
        \centering\sffamily\color{AlanyaGray}%
        \textbf{[ Capture a inserer ]}\\[10pt]\small #2\\[8pt]\texttt{captures/#1}%
      \end{minipage}}}%
    \caption{#2}%
  \end{figure}%
}

% Encadres typés
\newenvironment{note}[1][Note]{%
  \begin{tcolorbox}[enhanced, frame hidden, interior hidden,
    borderline west={3pt}{0pt}{AlanyaIndigo}, left=10pt, top=3pt, bottom=3pt,
    sharp corners, colback=white, before skip=\medskipamount, after skip=\medskipamount]%
  \small\textbf{\textcolor{AlanyaIndigoDark}{#1.}}\space\ignorespaces
}{\end{tcolorbox}}

\newenvironment{astuce}{%
  \begin{tcolorbox}[enhanced, frame hidden, interior hidden,
    borderline west={3pt}{0pt}{AlanyaGreen}, left=10pt, top=3pt, bottom=3pt,
    sharp corners, colback=white, before skip=\medskipamount, after skip=\medskipamount]%
  \small\textbf{\textcolor{AlanyaGreen}{Astuce.}}\space\ignorespaces
}{\end{tcolorbox}}

\newenvironment{attention}{%
  \begin{tcolorbox}[enhanced, frame hidden, interior hidden,
    borderline west={3pt}{0pt}{AlanyaAmber}, left=10pt, top=3pt, bottom=3pt,
    sharp corners, colback=white, before skip=\medskipamount, after skip=\medskipamount]%
  \small\textbf{\textcolor{AlanyaAmber}{Attention.}}\space\ignorespaces
}{\end{tcolorbox}}

\renewcommand{\arraystretch}{1.25}
\setlength{\parskip}{4pt}
\setlength{\parindent}{0pt}

% --------------------------------------------------------------------
\begin{document}

% ====================================================================
%  PAGE DE GARDE
% ====================================================================
\begin{titlepage}
\thispagestyle{empty}
\begin{tikzpicture}[remember picture, overlay]
  \fill[AlanyaIndigo]
    (current page.north west) rectangle ([yshift=-7.8cm]current page.north east);
  \fill[AlanyaIndigoDark]
    ([yshift=-7.8cm]current page.north west) rectangle ([yshift=-7.2cm]current page.north east);
\end{tikzpicture}

% -- Zone superieure : logos + nom de l'ecole --
\vspace*{0.45cm}
{\color{white}\sffamily
\begin{minipage}[c]{0.38\linewidth}
  \centering
  \IfFileExists{images/logo-enspy.png}{%
    \includegraphics[width=3.2cm]{images/logo-enspy.png}}{%
    \fbox{\begin{minipage}[c][2.8cm][c]{3.4cm}
      \centering\footnotesize\color{AlanyaGray}
      \textbf{Logo ENSPY}\\[3pt]\texttt{images/logo-enspy.png}
    \end{minipage}}}
\end{minipage}%
\hfill%
\begin{minipage}[c]{0.58\linewidth}
  \raggedleft
  {\small UNIVERSIT\'{E} DE YAOUND\'{E}~I}\\[4pt]
  {\small\bfseries \'{E}COLE NATIONALE SUP\'{E}RIEURE}\\
  {\small\bfseries POLYTECHNIQUE DE YAOUND\'{E}}\\[5pt]
  {\footnotesize D\'{e}partement G\'{e}nie Informatique}\\[3pt]
  {\footnotesize\itshape\color{AlanyaContainer} Projet Base de Donn\'{e}es}
\end{minipage}}

\vspace{1.9cm}

% -- Corps central --
\begin{center}\sffamily
  {\small\bfseries PROJET BASE DE DONN\'{E}ES}\\[10pt]
  {\color{AlanyaLine}\rule{0.6\linewidth}{0.6pt}}\\[22pt]
  {\fontsize{34}{38}\selectfont\bfseries\color{AlanyaIndigoDark} Manuel d'utilisation}\\[8pt]
  {\large\color{AlanyaGray} Application mobile de communication Alanya}\\[6pt]
  {\normalsize\itshape\color{AlanyaGray}
   Messagerie \textbullet{} Appels \textbullet{} R\'{e}unions \textbullet{} Statuts}\\[22pt]
  {\color{AlanyaLine}\rule{0.6\linewidth}{0.6pt}}\\[18pt]
  \IfFileExists{images/logo-Alanya.png}{%
    \includegraphics[width=5cm]{images/logo-Alanya.png}}{%
    \fbox{\begin{minipage}[c][3.6cm][c]{5.5cm}
      \centering\sffamily\color{AlanyaGray}
      \textbf{Logo Alanya}\\[3pt]\texttt{images/logo-Alanya.png}
    \end{minipage}}}
\end{center}

\vfill

% -- Informations academiques --
\begin{center}\sffamily
  \begin{tabular}{@{}r@{\quad}l@{}}
    \thd{Groupe N\textsuperscript{o}}        & 1\\[3pt]
    \thd{Sous la supervision de}             & Dr.\ NANA MBINKEU\\[3pt]
    \thd{Ann\'{e}e acad\'{e}mique}           & 2025-2026\\[3pt]
    \thd{Date}                               & \today\\
  \end{tabular}
\end{center}

\vspace{0.6cm}
\begin{center}\color{white}\small\sffamily
  Guide destin\'{e} aux utilisateurs de l'application
\end{center}
\end{titlepage}

% ====================================================================
%  TABLE DES MATIERES
% ====================================================================
\renewcommand{\contentsname}{Table des mati\`{e}res}
\tableofcontents
\listoffigures

% ====================================================================
\chapter{Introduction}
\label{chap:intro}

\section{\`{A} propos d'Alanya}
\textbf{Alanya} est une application mobile de communication qui r\'{e}unit, en un seul
endroit, la \textbf{messagerie} instantan\'{e}e, les \textbf{appels} audio et vid\'{e}o, les
\textbf{r\'{e}unions} en ligne et les \textbf{statuts} \'{e}ph\'{e}m\`{e}res. Elle vous permet
d'\'{e}changer des messages (texte, photos, vid\'{e}os, fichiers, messages vocaux),
de passer des appels individuels ou de groupe, d'organiser des visioconf\'{e}rences
planifi\'{e}es et de partager des statuts visibles pendant 24~heures.

\section{\`{A} qui s'adresse ce manuel}
Ce guide s'adresse \`{a} \textbf{toute personne utilisant Alanya} au quotidien. Aucune
connaissance technique n'est requise~: chaque fonctionnalit\'{e} est pr\'{e}sent\'{e}e
\textbf{pas \`{a} pas}, accompagn\'{e}e de captures d'\'{e}cran. Un chapitre distinct
(chapitre~\ref{chap:admin}) est r\'{e}serv\'{e} aux \textbf{administrateurs} de la plateforme.

\section{Conventions du document}
\begin{itemize}[leftmargin=1.3em]
  \item \bouton{Boutons} et commandes \`{a} toucher apparaissent en \textbf{gras sans serif}~;
  \item les noms de champs et d'onglets sont indiqu\'{e}s entre guillemets, par ex.\ \champ{Mot de passe}~;
  \item les proc\'{e}dures num\'{e}rot\'{e}es se lisent dans l'ordre indiqu\'{e}.
\end{itemize}
Trois types d'encadr\'{e}s ponctuent le texte~:
\begin{note}
information utile ou pr\'{e}cision de contexte.
\end{note}
\begin{astuce}
conseil pratique pour gagner du temps.
\end{astuce}
\begin{attention}
point de vigilance \`{a} ne pas n\'{e}gliger.
\end{attention}

\section{Pr\'{e}requis}
\begin{itemize}[leftmargin=1.3em]
  \item Un t\'{e}l\'{e}phone \textbf{Android} compatible.
  \item Une \textbf{connexion Internet} (Wi-Fi ou donn\'{e}es mobiles). Certaines actions
        (envoyer un message, lire l'historique) restent disponibles \textbf{hors ligne}
        et se synchronisent au retour du r\'{e}seau~; d'autres (nouvel appel, nouvelle
        discussion) n\'{e}cessitent une connexion active.
  \item Un \textbf{compte Alanya} (voir chapitre~\ref{chap:premierspas}).
\end{itemize}

% ====================================================================
\chapter{Premiers pas}
\label{chap:premierspas}

\section{Installer l'application}
Installez Alanya sur votre t\'{e}l\'{e}phone \`{a} partir du fichier d'installation (APK)
qui vous a \'{e}t\'{e} fourni, ou depuis la boutique d'applications. Ouvrez ensuite
l'application~: le premier \'{e}cran affich\'{e} est l'\'{e}cran de \textbf{connexion}.

\section{Cr\'{e}er un compte (inscription)}
Si vous n'avez pas encore de compte~:
\begin{enumerate}[leftmargin=1.5em]
  \item Sur l'\'{e}cran de connexion, touchez \bouton{S'inscrire}.
  \item Renseignez les champs \champ{Nom complet}, \champ{Pseudo},
        \champ{Adresse e-mail} et \champ{Mot de passe} (6\,caract\`{e}res minimum).
        L'ic\^{o}ne en forme d'\oe il permet d'afficher ou masquer le mot de passe.
  \item Touchez \bouton{Inscription} pour cr\'{e}er le compte.
  \item Une fen\^{e}tre affiche alors votre \textbf{identifiant Alanya}
        (\champ{Votre identifiant Talky})~: un num\'{e}ro \`{a} \textbf{6\,chiffres}.
        \textbf{Notez-le pr\'{e}cieusement}~: c'est lui qui vous servira \`{a} vous
        connecter et que vos contacts utiliseront pour vous joindre.
\end{enumerate}

\begin{attention}
L'identifiant \`{a} 6\,chiffres remplace le num\'{e}ro de t\'{e}l\'{e}phone classique pour la
connexion. Si vous l'oubliez, vous ne pourrez plus vous connecter~: conservez-le.
\end{attention}

\capture{01-signup.png}{\'{E}cran d'inscription et fen\^{e}tre affichant l'identifiant Alanya.}

\section{Se connecter}
\begin{enumerate}[leftmargin=1.5em]
  \item Dans le champ \champ{T\'{e}l\'{e}phone Alanya}, saisissez votre identifiant \`{a}
        6\,chiffres (par exemple \texttt{340364}).
  \item Saisissez votre \champ{Mot de passe}.
  \item Touchez \bouton{Se connecter}.
\end{enumerate}
En cas d'erreur (identifiant ou mot de passe incorrect), un message s'affiche en
rouge sous le formulaire.

\capture{02-login.png}{\'{E}cran de connexion.}

\section{Mot de passe oubli\'{e}}
Si vous avez oubli\'{e} votre mot de passe, touchez \bouton{Mot de passe oubli\'{e} ?}
sur l'\'{e}cran de connexion, puis suivez les \textbf{trois \'{e}tapes}~:
\begin{enumerate}[leftmargin=1.5em]
  \item \textbf{Adresse e-mail}~: saisissez l'e-mail de votre compte et touchez
        \bouton{Envoyer le code}. Un code \`{a} 6\,chiffres (OTP) vous est envoy\'{e} par e-mail.
  \item \textbf{V\'{e}rification}~: saisissez le code re\c{c}u, puis validez. Si vous ne
        l'avez pas re\c{c}u, touchez \bouton{Renvoyer le code}.
  \item \textbf{Nouveau mot de passe}~: saisissez le \champ{Nouveau mot de passe} et
        sa confirmation, puis validez. Le message \champ{Mot de passe r\'{e}initialis\'{e}
        avec succ\`{e}s} confirme l'op\'{e}ration et vous ramène \`{a} la connexion.
\end{enumerate}

\begin{note}
Le code OTP a une dur\'{e}e de validit\'{e} limit\'{e}e (10\,minutes). Pens\'{e} \`{a}
v\'{e}rifier votre dossier \og{}courrier ind\'{e}sirable\fg{} si vous ne le voyez pas.
\end{note}

\capture{03-otp.png}{R\'{e}initialisation du mot de passe par code OTP.}

\section{Autorisations demand\'{e}es}
Lors de la premi\`{e}re utilisation de certaines fonctions, Android vous demande
d'autoriser l'application. Acceptez pour que la fonction concern\'{e}e marche~:
\begin{table}[H]
\centering\small
\begin{tabularx}{\linewidth}{@{}l X@{}}
\toprule
\thd{Autorisation} & \thd{Demand\'{e}e quand\ldots}\\
\midrule
Micro    & vous passez un appel, rejoignez une r\'{e}union ou enregistrez un message vocal.\\
Cam\'{e}ra & vous passez un appel vid\'{e}o, rejoignez une r\'{e}union vid\'{e}o, ou prenez une photo (statut, profil).\\
Notifications & au d\'{e}marrage, pour vous pr\'{e}venir des messages, appels et r\'{e}unions.\\
Fichiers / Galerie & vous envoyez une photo, une vid\'{e}o ou un document.\\
\bottomrule
\end{tabularx}
\caption{Autorisations et moment o\`{u} elles sont demand\'{e}es.}
\end{table}

\begin{attention}
Si vous refusez le micro ou la cam\'{e}ra, les appels et r\'{e}unions ne pourront pas
fonctionner correctement. Vous pourrez r\'{e}activer ces autorisations dans les
r\'{e}glages Android de l'application.
\end{attention}

\section{D\'{e}couvrir l'\'{e}cran d'accueil}
Une fois connect\'{e}, vous arrivez sur l'\'{e}cran d'accueil. Une \textbf{barre de
navigation} en bas de l'\'{e}cran donne acc\`{e}s aux cinq sections de l'application.
Touchez une ic\^{o}ne pour changer de section.

\begin{figure}[H]
\centering
\begin{tikzpicture}[font=\sffamily\footnotesize]
  \node[draw=AlanyaLine, rounded corners=8pt, fill=boxBg, minimum width=15cm,
        minimum height=1.5cm] (bar) {};
  \foreach \i/\lab [evaluate=\i as \x using {-6+\i*3}] in
    {0/Discussions, 1/Appels, 2/Statuts, 3/R\'{e}unions, 4/Profil} {
    \node[text=AlanyaIndigo] at (\x,0.28) {$\bullet$};
    \node[text=AlanyaInk]    at (\x,-0.28) {\lab};
  }
\end{tikzpicture}
\caption{Barre de navigation : les cinq sections de l'application.}
\label{fig:nav}
\end{figure}

\begin{table}[H]
\centering\small
\begin{tabularx}{\linewidth}{@{}l X@{}}
\toprule
\thd{Section} & \thd{R\^{o}le}\\
\midrule
\bouton{Discussions} & Vos conversations et messages (chapitre~\ref{chap:chats}).\\
\bouton{Appels}      & Historique et lancement d'appels (chapitre~\ref{chap:appels}).\\
\bouton{Statuts}     & Statuts \'{e}ph\'{e}m\`{e}res (chapitre~\ref{chap:statuts}).\\
\bouton{R\'{e}unions} & R\'{e}unions planifi\'{e}es (chapitre~\ref{chap:reunions}).\\
\bouton{Profil}      & Votre profil et les r\'{e}glages (chapitre~\ref{chap:profil}).\\
\bottomrule
\end{tabularx}
\caption{Les cinq sections de la barre de navigation.}
\end{table}

% ====================================================================
\chapter{Messagerie (Discussions)}
\label{chap:chats}

\section{La liste des discussions}
La section \bouton{Discussions} affiche toutes vos conversations, la plus
r\'{e}cente en haut. Chaque ligne montre la photo du contact, son nom, un aper\c{c}u du
dernier message et l'heure. Un \textbf{point vert} sur la photo indique que le
contact est en ligne~; un \textbf{badge} indique le nombre de messages non lus.

\subsection{Rechercher et filtrer}
\begin{itemize}[leftmargin=1.3em]
  \item Touchez l'ic\^{o}ne \textbf{loupe} pour rechercher une conversation par nom ou pseudo.
  \item Les filtres permettent d'afficher~: \champ{Tous}, \champ{Discussions} (\'{e}changes
        individuels), \champ{Groupes}, \champ{Non lus} ou \champ{Archiv\'{e}s}.
\end{itemize}

\subsection{Actions rapides sur une conversation}
Faites un \textbf{appui long} sur une conversation pour ouvrir le menu d'actions~:
\begin{itemize}[leftmargin=1.3em]
  \item \bouton{\'{E}pingler} / \bouton{D\'{e}s\'{e}pingler}~: la conversation reste en haut de la liste~;
  \item \bouton{Archiver} / \bouton{D\'{e}sarchiver}~: la range dans \champ{Archiv\'{e}s}~;
  \item \bouton{Supprimer}~: retire la conversation de votre liste.
\end{itemize}

\capture{04-chats-liste.png}{Liste des discussions avec barre de recherche et filtres.}

\section{D\'{e}marrer une nouvelle discussion}
Touchez le bouton \textbf{+} (en bas \`{a} droite), choisissez un utilisateur, puis
commencez \`{a} \'{e}changer.

\begin{note}
Le bouton \textbf{+} est indisponible hors ligne~: un message vous le signale.
Reconnectez-vous pour d\'{e}marrer une nouvelle discussion.
\end{note}

\section{Cr\'{e}er un groupe}
\begin{enumerate}[leftmargin=1.5em]
  \item Touchez \bouton{Cr\'{e}er un groupe}.
  \item S\'{e}lectionnez les membres \`{a} ajouter.
  \item Donnez un \textbf{nom} (et \'{e}ventuellement une \textbf{photo}) au groupe.
  \item Touchez \bouton{Cr\'{e}er le groupe}.
\end{enumerate}
Dans un groupe existant, ouvrez la fiche du groupe (en touchant son nom) pour voir
les membres~; l'administrateur du groupe peut ajouter des membres
(\bouton{Ajouter des participants}), renommer le groupe ou changer sa photo.

\section{Envoyer un message}
Ouvrez une conversation, puis~:
\begin{itemize}[leftmargin=1.3em]
  \item \textbf{Texte}~: \'{e}crivez dans le champ de saisie et touchez le bouton d'envoi.
  \item \textbf{\'{E}moji}~: touchez l'ic\^{o}ne \'{e}moji pour ouvrir le s\'{e}lecteur.
  \item \textbf{Pi\`{e}ce jointe}~: touchez l'ic\^{o}ne \textbf{trombone} puis choisissez
        \bouton{Galerie} (photo/vid\'{e}o existante), \bouton{Cam\'{e}ra} (prendre une photo),
        ou un \textbf{fichier} (document, max 50\,Mo).
  \item \textbf{Message vocal}~: \textbf{maintenez} l'ic\^{o}ne \textbf{micro} enfonc\'{e}e pour
        enregistrer (un minuteur s'affiche), puis rel\^{a}chez pour envoyer. Pour annuler,
        glissez le doigt avant de rel\^{a}cher.
\end{itemize}

\capture{05-chat-detail.png}{\'{E}cran de conversation et options d'envoi (pi\`{e}ce jointe, vocal).}

\section{Agir sur un message}
Faites un \textbf{appui long} sur un message pour~:
\begin{itemize}[leftmargin=1.3em]
  \item \bouton{R\'{e}pondre}~: cite le message dans votre r\'{e}ponse~;
  \item \bouton{Copier}~: copie le texte (messages texte uniquement)~;
  \item \bouton{Modifier}~: \'{e}dite l'un de \emph{vos} messages texte~;
  \item \bouton{Supprimer pour moi}~: le retire de votre affichage~;
  \item \bouton{Supprimer pour tous}~: le retire pour tout le monde (vos messages).
\end{itemize}
Un message supprim\'{e} appara\^{i}t comme \champ{Ce message a \'{e}t\'{e} supprim\'{e}}.

\section{Comprendre les accus\'{e}s et la saisie}
Sous votre dernier message envoy\'{e}, une ic\^{o}ne indique son \'{e}tat~:
\begin{table}[H]
\centering\small
\begin{tabularx}{\linewidth}{@{}c l X@{}}
\toprule
\thd{Ic\^{o}ne} & \thd{\'{E}tat} & \thd{Signification}\\
\midrule
\textcolor{AlanyaGray}{$\circlearrowright$} & En attente & Le message est en cours d'envoi.\\
$\checkmark$ & Envoy\'{e} & Le serveur a re\c{c}u le message.\\
$\checkmark\checkmark$ \textcolor{AlanyaGray}{(gris)} & Distribu\'{e} & Le message est arriv\'{e} sur l'appareil du destinataire.\\
$\checkmark\checkmark$ \textcolor{AlanyaBlue}{(bleu)} & Lu & Le destinataire a ouvert le message.\\
\textcolor{red}{\textbf{!}} & \'{E}chec & L'envoi a \'{e}chou\'{e}~; touchez le message pour r\'{e}essayer.\\
\bottomrule
\end{tabularx}
\caption{Ic\^{o}nes d'accus\'{e} de message.}
\label{tab:accuses}
\end{table}
Lorsque votre interlocuteur \'{e}crit, la mention \champ{\ldots est en train d'\'{e}crire}
appara\^{i}t en haut de la conversation.

% ====================================================================
\chapter{Appels}
\label{chap:appels}

\section{L'historique des appels}
La section \bouton{Appels} liste vos appels r\'{e}cents avec, pour chacun~: la photo et
le nom du contact, le type (\champ{Appel vocal} ou \champ{Appel vid\'{e}o}), le sens
(\'{e}mis ou re\c{c}u), la dur\'{e}e et la date. La \textbf{loupe} permet de rechercher
dans l'historique.

\section{Passer un appel}
Vous pouvez lancer un appel de trois fa\c{c}ons~:
\begin{itemize}[leftmargin=1.3em]
  \item \textbf{Depuis une discussion}~: touchez l'ic\^{o}ne \textbf{combin\'{e}} (appel vocal)
        ou \textbf{cam\'{e}ra} (appel vid\'{e}o) en haut de la conversation~;
  \item \textbf{Depuis le clavier}~: dans \bouton{Appels}, ouvrez le clavier num\'{e}rique,
        composez l'identifiant Alanya \`{a} 6\,chiffres (la fiche du contact appara\^{i}t
        s'il existe), puis touchez le bouton d'appel~;
  \item \textbf{Depuis un contact}~: ouvrez la fiche d'un contact et choisissez l'appel.
\end{itemize}

\capture{06-keypad.png}{Clavier de num\'{e}rotation par identifiant Alanya.}

\section{Recevoir un appel}
Lorsqu'un appel arrive, un \'{e}cran plein affiche le nom et la photo de l'appelant
ainsi que le type d'appel~:
\begin{itemize}[leftmargin=1.3em]
  \item \bouton{Accepter} (bouton vert)~: d\'{e}marre l'appel~;
  \item \bouton{Refuser} (bouton rouge)~: rejette l'appel.
\end{itemize}

\section{Pendant l'appel}
L'\'{e}cran d'appel affiche la dur\'{e}e \'{e}coul\'{e}e et l'\'{e}tat (\champ{Connexion\ldots}
puis \champ{Connect\'{e}}). Les commandes disponibles~:
\begin{table}[H]
\centering\small
\begin{tabularx}{\linewidth}{@{}l X@{}}
\toprule
\thd{Commande} & \thd{Effet}\\
\midrule
Micro       & Coupe ou r\'{e}active votre microphone.\\
Haut-parleur & Bascule entre \'{e}couteur et haut-parleur.\\
Cam\'{e}ra   & Active ou coupe votre vid\'{e}o (appels vid\'{e}o).\\
Changer de cam\'{e}ra & Bascule entre cam\'{e}ra avant et arri\`{e}re.\\
Raccrocher (rouge) & Termine l'appel.\\
\bottomrule
\end{tabularx}
\caption{Commandes pendant un appel.}
\end{table}

\capture{07-appel.png}{\'{E}cran d'appel en cours et ses commandes.}

\section{Appels de groupe}
Alanya permet les appels \`{a} plusieurs (jusqu'\`{a} 9\,participants). Les vid\'{e}os des
participants s'affichent en mosa\"{i}que~; les commandes (micro, haut-parleur,
cam\'{e}ra, raccrocher) sont identiques \`{a} celles d'un appel individuel. Tant que
personne n'a rejoint, l'\'{e}cran indique \champ{En attente des participants\ldots}.

% ====================================================================
\chapter{R\'{e}unions}
\label{chap:reunions}

\section{La liste des r\'{e}unions}
La section \bouton{R\'{e}unions} organise vos r\'{e}unions en trois onglets~:
\champ{\`{A} venir}, \champ{En cours} et \champ{Pass\'{e}es}. Chaque carte indique le
titre, la date et l'heure, l'organisateur (\champ{H\^{o}te}) et le nombre de
participants. La \textbf{loupe} permet de filtrer par titre.

\section{Planifier une r\'{e}union}
\begin{enumerate}[leftmargin=1.5em]
  \item Touchez le bouton \textbf{+} pour ouvrir l'\'{e}cran de planification.
  \item Saisissez le \textbf{titre}, la \textbf{date} et l'\textbf{heure}.
  \item S\'{e}lectionnez les \textbf{participants} \`{a} inviter.
  \item Choisissez le type (audio ou vid\'{e}o), puis validez.
\end{enumerate}

\section{Inviter, accepter, refuser}
Depuis le \champ{D\'{e}tail de la r\'{e}union}, l'organisateur peut toucher
\bouton{Ajouter des participants} pour inviter d'autres personnes. Les invit\'{e}s
re\c{c}oivent une notification et peuvent \textbf{accepter} ou \textbf{refuser}
l'invitation. L'organisateur peut aussi \bouton{Annuler la r\'{e}union}.

\capture{08-meeting-detail.png}{D\'{e}tail d'une r\'{e}union et liste des participants.}

\section{La salle d'attente}
Avant d'entrer dans une r\'{e}union, vous passez par une \textbf{salle d'attente} qui
affiche un aper\c{c}u de votre cam\'{e}ra. Vous pouvez y \textbf{activer ou couper} votre
micro et votre cam\'{e}ra (\champ{Cam\'{e}ra active} / \champ{Cam\'{e}ra coup\'{e}e}) avant de
toucher \bouton{Rejoindre}.

\section{Participer \`{a} la r\'{e}union}
Une fois entr\'{e}, les vid\'{e}os des participants s'affichent en mosa\"{i}que (votre
propre vignette est rep\'{e}r\'{e}e). Les commandes permettent de couper le micro,
d'activer/couper la cam\'{e}ra, d'ouvrir le \bouton{Chat} de la r\'{e}union et de
\bouton{Quitter}.

\capture{09-meeting-live.png}{R\'{e}union en cours (mosa\"{i}que des participants).}

% ====================================================================
\chapter{Statuts}
\label{chap:statuts}

\section{Voir les statuts}
La section \bouton{Statuts} pr\'{e}sente d'abord \champ{Mon statut}, puis les statuts
de vos contacts r\'{e}partis en \champ{R\'{e}cents} (non encore vus) et \champ{D\'{e}j\`{a}
vus}. Touchez le statut d'un contact pour l'ouvrir en plein \'{e}cran~:
\begin{itemize}[leftmargin=1.3em]
  \item \textbf{glissez} vers la gauche/droite pour passer au statut suivant/pr\'{e}c\'{e}dent~;
  \item \textbf{touchez} pour mettre en pause, touchez \`{a} nouveau pour reprendre.
\end{itemize}

\section{Cr\'{e}er un statut}
Touchez \champ{Appuyer pour ajouter votre statut} (ou le bouton de cr\'{e}ation), puis
choisissez le type~:
\begin{table}[H]
\centering\small
\begin{tabularx}{\linewidth}{@{}l X@{}}
\toprule
\thd{Type} & \thd{\'{E}tapes}\\
\midrule
Texte & \'{E}crivez votre texte, choisissez une \champ{Couleur de fond}, puis touchez \bouton{Publier}.\\
Photo & Touchez \bouton{Ajouter une photo} (\bouton{Galerie} ou \bouton{Appareil}), ajoutez une description, puis \bouton{Publier}.\\
Vid\'{e}o & Touchez \bouton{Ajouter une vid\'{e}o}, ajoutez une description, puis \bouton{Publier}.\\
Audio & Enregistrez un \champ{Message vocal} (possibilit\'{e} de \bouton{R\'{e}enregistrer}), puis \bouton{Publier}.\\
\bottomrule
\end{tabularx}
\caption{Les quatre types de statut.}
\end{table}

\begin{note}
Un statut reste visible \textbf{24\,heures}, puis dispara\^{i}t automatiquement.
\end{note}

\capture{10-status-create.png}{Cr\'{e}ation d'un statut (texte avec couleur de fond).}

\section{Qui a vu mon statut~? Aimer un statut}
Ouvrez \champ{Mon statut} et touchez le nombre de vues pour voir \textbf{qui} a
regard\'{e} votre statut et \textbf{quand} (\champ{Aucune vue} s'il n'y en a pas
encore). En regardant le statut d'un contact, touchez l'ic\^{o}ne \textbf{c\oe ur} pour
l'aimer, ou \champ{R\'{e}pondre au statut\ldots} pour envoyer une r\'{e}ponse en message
priv\'{e}.

% ====================================================================
\chapter{Contacts}
\label{chap:contacts}

\section{Rechercher et ajouter un contact}
Touchez \bouton{Nouveau contact} puis recherchez un utilisateur (par identifiant
ou par nom). Vous pouvez l'ajouter \`{a} vos \textbf{contacts pr\'{e}f\'{e}r\'{e}s} pour le
retrouver rapidement. \bouton{Nouveau groupe} permet de cr\'{e}er directement un groupe.

\section{Contacts pr\'{e}f\'{e}r\'{e}s}
Les \champ{Contacts pr\'{e}f\'{e}r\'{e}s} sont accessibles rapidement (notamment depuis votre
profil et l'\'{e}cran d'appel). Ajoutez ou retirez un favori depuis la fiche du contact.

\section{Bloquer un contact}
Depuis la fiche d'un contact, touchez \bouton{Bloquer le contact} et confirmez
(\champ{Bloquer ce contact ?}). Un contact bloqu\'{e} ne peut plus vous joindre. Pour
annuler, touchez \bouton{D\'{e}bloquer}.

\capture{11-contact.png}{Fiche d'un contact (appeler, message, favori, bloquer).}

% ====================================================================
\chapter{Profil et r\'{e}glages}
\label{chap:profil}

\section{Votre profil}
La section \bouton{Profil} affiche votre photo, votre nom, votre pseudo, votre
identifiant Alanya et vos contacts pr\'{e}f\'{e}r\'{e}s. De l\`{a}, vous pouvez ouvrir
\bouton{Modifier le profil}, les \bouton{Param\`{e}tres}, l'espace \bouton{Admin}
(si votre compte est administrateur) et vous \bouton{D\'{e}connecter}.

\section{Modifier le profil}
Dans \bouton{Modifier le profil}~:
\begin{itemize}[leftmargin=1.3em]
  \item touchez votre photo pour \bouton{Prendre une photo}, \bouton{Choisir depuis la
        galerie} ou \bouton{Supprimer la photo}~;
  \item modifiez votre \champ{Nom}, votre \champ{Pseudo}, votre e-mail, votre pays~;
  \item touchez \bouton{Enregistrer}. Le message \champ{Profil mis \`{a} jour} confirme.
\end{itemize}

\capture{12-edit-profile.png}{\'{E}cran de modification du profil.}

\section{R\'{e}glages}
L'\'{e}cran \bouton{Param\`{e}tres} regroupe~:
\begin{itemize}[leftmargin=1.3em]
  \item \champ{Apparence}~: th\`{e}me \champ{Clair} ou \champ{Sombre}~;
  \item \champ{Notifications}~: sonneries de message, de groupe et d'appel~;
  \item \champ{Confidentialit\'{e}}~: blocage de contacts, messages \'{e}ph\'{e}m\`{e}res~;
  \item \champ{Compte} / \champ{S\'{e}curit\'{e}}~: gestion du compte~;
  \item \champ{Stockage et donn\'{e}es}.
\end{itemize}

\section{Se d\'{e}connecter}
Touchez \bouton{D\'{e}connexion} dans votre profil. Vous revenez \`{a} l'\'{e}cran de
connexion. Vos donn\'{e}es restent associ\'{e}es \`{a} votre compte et seront restaur\'{e}es
\`{a} la prochaine connexion.

% ====================================================================
\chapter{Espace administrateur}
\label{chap:admin}

\begin{note}
Ce chapitre concerne uniquement les comptes \textbf{administrateur}. L'entr\'{e}e
\bouton{Admin} appara\^{i}t alors dans la section \bouton{Profil}.
\end{note}

\section{Le tableau de bord}
L'\'{e}cran \champ{Tableau de bord} pr\'{e}sente des statistiques globales sous forme de
cartes~: nombre total d'utilisateurs, utilisateurs \champ{En ligne},
\champ{Messages (7j)}, \champ{Appels (7j)}, \champ{Statuts (7j)} et nombre de
comptes \champ{Bannis}.

\capture{13-admin-dashboard.png}{Tableau de bord d'administration.}

\section{Rechercher et filtrer les utilisateurs}
Sous les statistiques, la liste des utilisateurs peut \^{e}tre~:
\begin{itemize}[leftmargin=1.3em]
  \item \textbf{filtr\'{e}e} par \champ{Tous}, \champ{En ligne}, \champ{Bannis} ou \champ{Admins}~;
  \item \textbf{recherch\'{e}e} par nom, pseudo, e-mail ou t\'{e}l\'{e}phone~;
  \item \textbf{parcourue} page par page lorsqu'il y a beaucoup de comptes.
\end{itemize}

\section{G\'{e}rer un utilisateur}
Touchez une ligne pour ouvrir la fiche d\'{e}taill\'{e}e d'un utilisateur~: profil,
activit\'{e} (messages, appels, statuts) et \champ{Connexions r\'{e}centes}. Les actions
disponibles~:
\begin{table}[H]
\centering\small
\begin{tabularx}{\linewidth}{@{}l X@{}}
\toprule
\thd{Action} & \thd{Effet}\\
\midrule
\bouton{Bannir} & Emp\^{e}che l'utilisateur de se connecter (un motif peut \^{e}tre saisi).\\
\bouton{D\'{e}bannir} & Lève le bannissement.\\
\bouton{Rendre admin} / \bouton{R\'{e}trograder} & Change le r\^{o}le du compte.\\
\bouton{Supprimer} & Supprime le compte (action irr\'{e}versible).\\
\bottomrule
\end{tabularx}
\caption{Actions d'administration sur un utilisateur.}
\end{table}

\begin{attention}
Le bannissement et la suppression ont un impact imm\'{e}diat sur l'utilisateur
concern\'{e}. La suppression est \textbf{irr\'{e}versible}~: v\'{e}rifiez avant de confirmer.
\end{attention}

% ====================================================================
\chapter{Foire aux questions \& d\'{e}pannage}
\label{chap:faq}

\begin{table}[H]
\centering\small
\begin{tabularx}{\linewidth}{@{}p{5.4cm} X@{}}
\toprule
\thd{Probl\`{e}me} & \thd{Que faire}\\
\midrule
Je ne re\c{c}ois pas le code OTP & V\'{e}rifiez votre dossier \og{}spam\fg{}, l'adresse e-mail saisie, puis touchez \bouton{Renvoyer le code} (le code expire au bout de 10\,minutes).\\
Connexion impossible & V\'{e}rifiez votre identifiant \`{a} 6\,chiffres et votre mot de passe~; en cas d'oubli, utilisez \bouton{Mot de passe oubli\'{e} ?}. Un compte banni ne peut pas se connecter.\\
Mon appel ne passe pas / pas de son ni d'image & Autorisez le \textbf{micro} et la \textbf{cam\'{e}ra} dans les r\'{e}glages Android, et v\'{e}rifiez votre connexion Internet.\\
Je ne re\c{c}ois pas les notifications & Autorisez les \textbf{notifications} pour Alanya dans les r\'{e}glages Android et activez-les dans \bouton{Param\`{e}tres} \textrightarrow{} \champ{Notifications}.\\
Mon message n'est pas envoy\'{e} (ic\^{o}ne d'erreur) & Vous \^{e}tes peut-\^{e}tre hors ligne~: le message part automatiquement au retour du r\'{e}seau. Vous pouvez aussi toucher le message pour r\'{e}essayer.\\
\bouton{Nouvelle discussion} indisponible & Cette action n\'{e}cessite une connexion~: reconnectez-vous.\\
\bottomrule
\end{tabularx}
\caption{Probl\`{e}mes courants et solutions.}
\end{table}

% ====================================================================
\appendix
\chapter{Annexe : glossaire et pictogrammes}

\section{Glossaire}
\begin{tabularx}{\linewidth}{@{}l X@{}}
\toprule
\thd{Terme} & \thd{D\'{e}finition}\\
\midrule
Identifiant Alanya & Num\'{e}ro \`{a} 6\,chiffres servant \`{a} se connecter et \`{a} \^{e}tre joint.\\
OTP & Code \`{a} usage unique re\c{c}u par e-mail pour r\'{e}initialiser le mot de passe.\\
Statut & Publication \'{e}ph\'{e}m\`{e}re visible 24\,heures.\\
Hors ligne & Sans connexion Internet~; certaines actions sont diff\'{e}r\'{e}es jusqu'au retour du r\'{e}seau.\\
Contact pr\'{e}f\'{e}r\'{e} & Contact mis en favori pour un acc\`{e}s rapide.\\
H\^{o}te & Organisateur d'une r\'{e}union.\\
\bottomrule
\end{tabularx}

\section{Ic\^{o}nes d'accus\'{e} de message}
Le rappel des ic\^{o}nes \champ{En attente}, \champ{Envoy\'{e}}, \champ{Distribu\'{e}},
\champ{Lu} et \champ{\'{E}chec} figure au tableau~\ref{tab:accuses}
(page~\pageref{tab:accuses}).

\end{document}
